LLM-based coding agents systematically declare tasks complete when work is
partial or incorrect, a phenomenon we term \emph{false completion}.
This failure mode arises from at least seven independent root causes spanning
model architecture, reinforcement-learning training methodology, and
tool-use infrastructure.
No prior work provides a unified taxonomy, quantitative evidence, or
mechanical mitigation for false completion in agentic LLM systems.
We present three contributions:
(1)~a taxonomy of 7~root causes, 10~failure patterns, and 16~failure modes
distributed across 6~lifecycle phases, synthesized from academic research,
vendor publications, and more than 20~community-reported incidents;
(2)~quantitative evidence including analysis of 6{,}852 sessions showing an
80$\times$ API cost increase and a 444\% rise in unverified edits following
infrastructure changes;
(3)~the Completion Integrity System, a three-layer mechanical verification
architecture that uses transcript analysis, file-system verification, and
skill-level gate markers to detect discrepancies between completion claims
and observable evidence.
Cross-model experimentation with Claude, Qwen, and Gemini on 20~standardized
tasks demonstrates [to be updated with experimental results].
Mechanical enforcement achieves [to be updated] reduction in false completions
compared to prompt-only mitigation.
These findings support what we call the \emph{Worker-Inspector Paradox}:
reliable verification of agentic LLM output requires external mechanical
checks independent of the model's own incentive structure.
